% ==========================================================
% 070_thresholds_and_auditability.tex
% Forecast Admissibility Surface (FAS)
% ==========================================================

\section{Thresholds, Determinism, and Auditability}
\label{sec:fas_thresholds_auditability}

Admissibility classification in the Forecast Admissibility Surface is governed by
a fixed set of threshold parameters applied to baseline error anatomy. These
thresholds define the boundary between \Allowed, \Conditional, and \Blocked\
regions of the slice domain. While the specific numerical values of thresholds
are domain-dependent, the manner in which they are applied is invariant.

\subsection{Threshold Parameterization}

FAS thresholds operate on three classes of observable signals: support size,
spike frequency, and tail magnitude. Support thresholds guard against
overinterpretation of sparse data, while spike and tail thresholds detect
structural instability in baseline error behavior. Each threshold is expressed
as a simple, interpretable condition (e.g., spike rate greater than or equal to a
specified level).

Thresholds are not learned, optimized, or adapted dynamically. They are selected
as governance parameters based on operational tolerance for instability and
risk. This design ensures that admissibility boundaries are explicit and
understandable, rather than emergent artifacts of model fitting.

\subsection{Deterministic Classification}

Given a fixed dataset, baseline forecast, slice definition, and threshold
configuration, FAS produces a unique admissibility surface. Classification rules
are evaluated independently for each slice, and the resulting admissibility
labels are fully determined by observed error anatomy.

This determinism is essential for operational governance. Admissibility decisions
must be reproducible across runs, environments, and stakeholders. Two analysts
applying FAS to the same inputs must arrive at the same classification, and any
change in outcome must be attributable to a documented change in data, slice
definition, or threshold configuration.

\subsection{Auditability and Traceability}

To support auditability, FAS associates each admissibility surface with a stable
representation of the threshold configuration used to produce it. Thresholds
are serialized in a canonical form and accompanied by a compact fingerprint that
uniquely identifies the configuration.

This fingerprint enables downstream artifacts—such as forecast panels,
evaluation reports, and governance decisions—to be traced back to the exact
admissibility criteria in force at the time of generation. When classifications
change, the cause can be unambiguously identified as a change in observed error
anatomy or a deliberate governance update to thresholds.

\subsection{Governance Implications}

By externalizing admissibility boundaries as explicit thresholds, FAS separates
structural decisions from modeling outcomes. Adjusting thresholds is a
governance action, not a modeling one. Such changes should be rare, deliberate,
and documented, reflecting shifts in operational risk tolerance or strategic
objectives rather than short-term performance fluctuations.

This separation reinforces the role of FAS as a governance primitive within the
Forecast Readiness Framework: a stable, auditable gate that constrains where
forecasting is permitted, independent of downstream optimization or evaluation
concerns.
